\documentclass[11pt]{article}
\usepackage[margin=0.9in]{geometry}
\usepackage{amsmath,amssymb,amsthm,mathtools}
\usepackage[T1]{fontenc}
\usepackage[utf8]{inputenc}
\usepackage{enumitem}
\usepackage{booktabs,longtable,array,seqsplit}
\usepackage{hyperref}
\usepackage{xcolor}
\usepackage{microtype}
\hypersetup{colorlinks=true,linkcolor=blue!45!black,citecolor=blue!45!black,urlcolor=blue!45!black}
\setlength{\emergencystretch}{6em}

\newtheorem{theorem}{Theorem}[section]
\newtheorem{proposition}[theorem]{Proposition}
\newtheorem{lemma}[theorem]{Lemma}
\newtheorem{corollary}[theorem]{Corollary}
\newtheorem{claim}[theorem]{Claim}
\theoremstyle{definition}
\newtheorem{definition}[theorem]{Definition}
\newtheorem{assumption}[theorem]{Assumption}
\theoremstyle{remark}
\newtheorem{remark}[theorem]{Remark}

\newcommand{\R}{\mathbb{R}}
\newcommand{\C}{\mathbb{C}}
\newcommand{\N}{\mathbb{N}}
\newcommand{\Hsp}{\mathcal{H}}
\newcommand{\Fsp}{\mathcal{F}}
\newcommand{\Id}{\mathrm{Id}}
\newcommand{\tr}{\operatorname{tr}}
\newcommand{\diag}{\operatorname{diag}}
\newcommand{\rank}{\operatorname{rank}}
\newcommand{\spec}{\operatorname{spec}}
\newcommand{\spanop}{\operatorname{span}}
\newcommand{\dist}{\operatorname{dist}}
\newcommand{\range}{\operatorname{ran}}
\newcommand{\Ker}{\operatorname{ker}}
\newcommand{\Pinch}{\mathcal{C}}
\newcommand{\Off}{\widetilde{\mathcal{C}}}
\newcommand{\Pperp}{P^{\perp}}
\newcommand{\Qperp}{Q^{\perp}}
\newcommand{\op}{\mathrm{op}}
\newcommand{\F}{\mathrm{F}}
\newcommand{\HS}{\mathrm{HS}}
\newcommand{\ip}[2]{\left\langle #1,#2\right\rangle}
\newcommand{\norm}[1]{\left\lVert #1\right\rVert}
\newcommand{\abs}[1]{\left\lvert #1\right\rvert}
\newcommand{\code}[1]{\nolinkurl{#1}}

\newenvironment{argument}{\paragraph{Argument.}\begin{enumerate}[leftmargin=2em]}{\end{enumerate}}
\newenvironment{proofoutline}{\paragraph{Proof architecture.}\begin{enumerate}[leftmargin=2em]}{\end{enumerate}}
\newcommand{\sourceanchor}[1]{\par\smallskip\noindent\textbf{Source anchor.} #1\par\smallskip}
\newcommand{\sourcecomment}[1]{\begin{remark}#1\end{remark}}
\newenvironment{sourceguide}{\begin{description}[style=nextline,leftmargin=3.3cm,labelwidth=3.1cm]}{\end{description}}
\newenvironment{formalizationmap}{\begin{longtable}{@{}p{0.20\textwidth}p{0.31\textwidth}p{0.40\textwidth}@{}}\toprule Source item & Modern mathematical content & Repository status \\\midrule\endfirsthead\toprule Source item & Modern mathematical content & Repository status \\\midrule\endhead}{\bottomrule\end{longtable}}

\title{Wedin (1972): residual perturbation bounds for singular subspaces}
\author{}
\date{Distilled reconstruction for the finite Davis--Kahan formalization}
\begin{document}
\maketitle

\section*{Source map and attribution boundary}
\begin{sourceguide}
\item[Primary source] Per-\AA ke Wedin, \emph{Perturbation Bounds in Connection with Singular Value Decomposition}, BIT 12 (1972), 99--111, DOI \href{https://doi.org/10.1007/BF01932678}{10.1007/BF01932678}.
\item[Source access] The publisher record, abstract, and archival scan were checked. The scan has imperfect OCR, so this note emphasizes the invariant theorem architecture and does not attach unstable equation numbers to every modernized formula.
\item[Paper scope] Perturbations of the invariant subspaces of $A^*A$ and $AA^*$; simultaneous left/right singular-subspace estimates; residual formulations; applications to least squares, pseudoinverses, and rank-deficient approximation.
\item[Formalization target] \code{DavisKahanTheory.SingularSubspace}, especially \code{rightSingularSubspace}, \code{leftSingularSubspace}, \code{hermitianDilation}, and the missing canonical Wedin wrapper over the existing finite spectral-subspace machinery.
\item[Warning] ``Wedin's theorem'' now names a family of variants. Gap orientation, inclusion of the rectangular zero singular value, residual convention, and norm scope must be stated explicitly. A casual formula without these choices is not source-faithful.
\end{sourceguide}

Wedin's central contribution is to treat the left and right singular spaces as a coupled perturbation problem. Applying an eigenvector theorem separately to $A^*A$ and $AA^*$ squares the spectral scale and replaces the perturbation by Gram differences. The residual formulation instead works on the singular-value scale and exposes the two equations that couple the left and right angles.

\section{Partitioned singular value decompositions}

Let $A:\Fsp\to\Hsp$ be a finite-dimensional operator. Choose a singular-value partition
\[
 A=U_1\Sigma_1V_1^*+U_0\Sigma_0V_0^*.
\]
Here $U_1,V_1$ span the selected left and right singular subspaces, while $U_0,V_0$ denote the complementary blocks. For a perturbed operator
\[
 \widetilde A=A+E,
\]
choose a corresponding partition
\[
 \widetilde A=\widetilde U_1\widetilde\Sigma_1\widetilde V_1^*
                 +\widetilde U_0\widetilde\Sigma_0\widetilde V_0^*.
\]
The selected dimensions are normally matched, but the ambient left and right dimensions may differ.

The directed principal-angle matrices are encoded by the cross blocks
\[
 U_0^*\widetilde U_1,
 \qquad
 V_0^*\widetilde V_1.
\]
Their singular values are the sines of the principal angles between the selected left subspaces and between the selected right subspaces, respectively. Thus a simultaneous Wedin conclusion controls both
\[
 \sin\Theta_U:=\sin\Theta(\range U_1,\range\widetilde U_1),
 \qquad
 \sin\Theta_V:=\sin\Theta(\range V_1,\range\widetilde V_1).
\]

\section{Residuals rather than Gram perturbations}

Define the residuals of the perturbed singular triplets relative to $A$ by
\begin{align}
 R&:=A\widetilde V_1-\widetilde U_1\widetilde\Sigma_1,\\
 S&:=A^*\widetilde U_1-\widetilde V_1\widetilde\Sigma_1.
\end{align}
Because $\widetilde A=A+E$ satisfies
\[
 \widetilde A\widetilde V_1=\widetilde U_1\widetilde\Sigma_1,
 \qquad
 \widetilde A^*\widetilde U_1=\widetilde V_1\widetilde\Sigma_1,
\]
one also has
\[
 R=-E\widetilde V_1,
 \qquad
 S=-E^*\widetilde U_1.
\]
Signs are immaterial for norm estimates, but a formal theorem should fix one convention and use it consistently.

Multiplying the residual equations by complementary singular vectors gives the coupled Sylvester system
\begin{align}
 \Sigma_0(V_0^*\widetilde V_1)
 -(U_0^*\widetilde U_1)\widetilde\Sigma_1
 &=U_0^*R,\label{W-coupled-1}\\
 \Sigma_0^*(U_0^*\widetilde U_1)
 -(V_0^*\widetilde V_1)\widetilde\Sigma_1
 &=V_0^*S.\label{W-coupled-2}
\end{align}
For the usual nonnegative diagonal singular-value blocks, $\Sigma_0^*=\Sigma_0$. The two unknown cross blocks are precisely the left and right sine matrices.

This is the source-faithful analytic kernel. It avoids the Gram identities
\[
 \widetilde A^*\widetilde A-A^*A=A^*E+E^*A+E^*E
\]
and
\[
 \widetilde A\widetilde A^*-AA^*=AE^*+EA^*+EE^*,
\]
whose norms scale with the largest singular value and whose spectral gaps are differences of squares.

\section{The gap and the rectangular zero mode}

A useful modern formulation introduces a selected perturbed singular set
\[
 \widetilde\Sigma_1=\diag(\widetilde\sigma_1,\ldots,\widetilde\sigma_r)
\]
and the complementary singular set of $A$. The required gap has the form
\[
 \delta
 \le
 \abs{\widetilde\sigma_i-\sigma_j}
 \quad\text{for every selected $i$ and complementary $j$},
\]
together with the corresponding sum separation generated by the signed dilation spectrum.

For a rectangular operator, the Hermitian dilation has an additional zero eigenspace. Consequently the complementary set must sometimes be enlarged by $0$:
\[
 \spec_{\mathrm{ext}}(\Sigma_0)
 =\spec(\Sigma_0)\cup\{0\}.
\]
The distance of the selected singular values from zero is not cosmetic. It prevents the selected positive dilation eigenspace from mixing with unmatched left or right null directions.

A robust finite hypothesis is therefore
\[
 \delta
 \le
 \dist\bigl(
   \spec(\widetilde\Sigma_1),
   \spec(\Sigma_0)\cup\{0\}_{\mathrm{ghost}}
 \bigr),
\]
with the zero adjoined exactly when dictated by the ambient dimensions and chosen blocks.

\section{The simultaneous sine estimate}

The characteristic Wedin estimate combines the two angle errors and the two residuals. In Frobenius form it has the schematic sharp-gap shape
\begin{equation}
 \left(
   \norm{\sin\Theta_U}_{\F}^2
   +\norm{\sin\Theta_V}_{\F}^2
 \right)^{1/2}
 \le
 \frac{
   \left(\norm{R}_{\F}^2+\norm{S}_{\F}^2\right)^{1/2}
 }{\delta}.
 \tag{W-F}
\end{equation}
The source develops more general norm bounds under the corresponding separation assumptions. Modern statements often package the pair as a block diagonal operator:
\[
 \diag(\sin\Theta_U,\sin\Theta_V)
\]
and compare it with
\[
 \diag(R,S).
\]
This is the cleanest route to Ky Fan and arbitrary unitarily invariant norms, because it preserves the pairing between left and right singular directions.

The proof is a two-equation analogue of a Sylvester estimate. Squaring and adding the scalar coordinate equations yields a denominator bounded below by the singular-value gap. Summing over singular directions gives the Frobenius theorem. For general unitarily invariant norms, one needs a majorization statement for the coupled block system rather than independent entrywise estimates.

\section{Hermitian dilation: the reusable Davis--Kahan route}

Define the Hermitian dilation
\[
 \mathcal H(A)
 =
 \begin{pmatrix}
 0&A\\
 A^*&0
 \end{pmatrix}
 \quad\text{on }\Hsp\oplus\Fsp.
\]
If $Av=\sigma u$ and $A^*u=\sigma v$, then
\[
 \mathcal H(A)\binom{u}{v}=\sigma\binom{u}{v},
 \qquad
 \mathcal H(A)\binom{u}{-v}=-\sigma\binom{u}{-v}.
\]
Thus the nonzero spectrum of $\mathcal H(A)$ is the signed singular spectrum of $A$.

The perturbation is
\[
 \mathcal H(\widetilde A)-\mathcal H(A)
 =\mathcal H(E),
\]
and in operator norm
\[
 \norm{\mathcal H(E)}_{\op}=\norm{E}_{\op}.
\]
A Davis--Kahan theorem for the selected positive signed singular cluster therefore gives a singular-subspace theorem after proving a geometric dictionary between:
\begin{itemize}[leftmargin=2em]
\item the spectral projector of the dilation;
\item the left and right singular projectors;
\item the principal angles in the direct sum; and
\item the desired separate or coupled angle norm.
\end{itemize}

The dilation route is highly reusable, but it is not automatically identical to the residual theorem. A perturbation-only estimate uses $\norm{E}$, whereas Wedin's residual form can be much smaller because it measures only $E\widetilde V_1$ and $E^*\widetilde U_1$. The canonical API should expose both layers rather than naming the easiest dilation corollary as the entire Wedin theorem.

\section{Projector and basis conclusions}

A subspace theorem controls projectors or principal angles, not a chosen singular-vector basis. If the selected singular values have multiplicity, individual singular vectors may rotate arbitrarily inside the selected subspace. Source-faithful conclusions should therefore be stated first as
\[
 \norm{P_{U_1}-P_{\widetilde U_1}},
 \qquad
 \norm{P_{V_1}-P_{\widetilde V_1}},
\]
or through the directed sine maps.

A basis-alignment conclusion requires an additional Procrustes choice. Under equal rank, there exist unitary matrices $Q_U,Q_V$ such that
\[
 \norm{U_1Q_U-\widetilde U_1}
 \quad\text{and}\quad
 \norm{V_1Q_V-\widetilde V_1}
\]
are controlled by the principal angles. This is downstream geometry, not part of the fundamental spectral separation theorem.

\section{Finite formalization plan}

\begin{formalizationmap}
Singular subspaces & Spectral subspaces of $A^*A$ and $AA^*$ for squared singular-value sets & \code{rightSingularSubspace} and \code{leftSingularSubspace} exist.\\
Hermitian dilation & Self-adjoint block operator with signed singular spectrum & \code{hermitianDilation}, \code{isSymmetric_hermitianDilation}, and \code{hermitianDilation_sq} exist.\\
Gram perturbation & Bounds for $\widetilde A^*\widetilde A-A^*A$ and the left analogue & Operator-norm Gram difference lemmas and separate sine bounds exist.\\
Dilation sine theorem & Apply finite Davis--Kahan to selected dilation spectral sets & \code{singularSubspace_dilation_sinTheta_le} is present, but the singular-angle dictionary remains the key wrapper seam.\\
Wedin residuals & Coupled equations \eqref{W-coupled-1}--\eqref{W-coupled-2} & Not yet packaged as a source-faithful residual theorem.\\
Rectangular zero mode & Enlarge the complementary singular set by zero where necessary & Must be made explicit in the gap predicate; do not bury it in dimension arithmetic.\\
Canonical Wedin API & Simultaneous left/right angle bound, then perturbation corollaries & Still missing; should distinguish residual and dilation variants.\\
\end{formalizationmap}

A productive theorem sequence is:
\begin{enumerate}[leftmargin=2em]
\item Prove the exact norm and adjoint identities for \code{hermitianDilation}.
\item Identify its positive and negative spectral subspaces with paired left/right singular subspaces.
\item Prove the principal-angle duplication formula for the paired direct-sum subspaces.
\item Expose a clean perturbation-only Wedin corollary from the existing Davis--Kahan theorem.
\item Separately formalize the coupled residual equations and derive the sharper residual theorem.
\item Add basis-alignment corollaries only after the projector-level theorem is stable.
\end{enumerate}

\section{Failure modes to exclude}

\begin{itemize}[leftmargin=2em]
\item Do not use only $\min\abs{\widetilde\sigma_i-\sigma_j}$ when the negative dilation branch or the rectangular zero eigenspace is also in the complement.
\item Do not identify a projector bound with a bound for individual singular vectors in the presence of multiplicity.
\item Do not silently replace residual norms by $\norm{E}$ and still call the result the full residual theorem.
\item Do not prove the left and right estimates independently if the intended theorem is a coupled Frobenius or UI-norm estimate.
\item Do not square the problem through Gram operators unless the resulting squared gap and perturbation dependence are explicitly desired.
\end{itemize}

\section*{Distilled conclusion}

Wedin's theorem is the singular-value counterpart of Davis--Kahan, but its natural primitive is a coupled pair of residual Sylvester equations. Hermitian dilation offers the shortest route to a reusable perturbation corollary; the residual system is the route to the source's sharper content. The finite Lean API should preserve this distinction, expose the zero-mode gap explicitly, and state conclusions at the subspace level before adding basis alignment.

\end{document}
